\chapter*{KẾT LUẬN VÀ HƯỚNG PHÁT TRIỂN}
\addcontentsline{toc}{chapter}{KẾT LUẬN VÀ HƯỚNG PHÁT TRIỂN}
\markboth{KẾT LUẬN VÀ HƯỚNG PHÁT TRIỂN}{}

\section*{1. Kết luận}
\addcontentsline{toc}{section}{1. Kết luận}

Luận văn đã nghiên cứu bài toán phân bổ tài nguyên trong hệ thống SISO STAR--RIS hỗ trợ
RSMA với một trạm gốc, bốn thiết bị người dùng và một bề mặt STAR--RIS thụ động. Trên cùng
một mô hình tính toán, nghiên cứu đồng thời xét năm nhóm biến điều khiển gồm công suất
luồng chung và các luồng riêng, tỷ lệ phân bổ tốc độ luồng chung, hệ số chia năng lượng
truyền và phản xạ, pha phía truyền và pha phía phản xạ. Mục tiêu của bài toán là nâng cao
tổng tốc độ của hệ thống trong khi vẫn duy trì yêu cầu chất lượng dịch vụ tối thiểu cho
từng người dùng.

Trước khi lựa chọn thuật toán học, luận văn tập trung phân tích ý nghĩa vật lý của cấu
trúc SISO--RSMA. Trong hệ SISO, trạm gốc không có véc-tơ tạo búp sóng (beamforming) nhiều
chiều để tách các luồng riêng theo không gian. Vì vậy, khi nhiều luồng riêng cùng được
phát trên một tài nguyên, mỗi người dùng phải xem toàn bộ các luồng riêng còn lại như
nhiễu. Ở miền mà công suất nhiễu lớn hơn đáng kể so với tạp âm, độ lợi kênh xuất hiện đồng
thời ở tử số và mẫu số của SINR luồng riêng, khiến lợi ích của việc tăng đồng loạt công
suất bị hạn chế. Kết quả phân tích này giải thích vì sao phân bổ công suất riêng tập trung
vào một luồng chiếm ưu thế thường tạo điều kiện thuận lợi hơn cho tổng tốc độ trong mô
hình đang xét.

Luồng chung của RSMA đồng thời có vai trò quan trọng đối với QoS. Tốc độ luồng chung bị
giới hạn bởi người dùng khó giải mã nhất, nhưng sau khi xác định, phần tốc độ này có thể
được phân bổ linh hoạt cho các người dùng. Nhờ đó, một người dùng nhận ít công suất luồng
riêng vẫn có thể đạt tốc độ tối thiểu bằng phần tốc độ luồng chung. Phân tích trong miền
tỷ số tín hiệu trên tạp âm cao còn cho thấy, dưới giả thiết một luồng riêng chiếm ưu thế
và khả năng giải mã luồng chung đủ tốt, phần tăng của tốc độ luồng chung có thể bù cho
phần giảm của tốc độ luồng riêng khi tỷ lệ công suất luồng chung tăng. Đây là một xấp xỉ
dùng để tạo trực giác thiết kế, không phải chứng minh rằng một tỷ lệ công suất cụ thể là
tối ưu cho mọi trạng thái kênh.

Từ những quan sát trên, luận văn xây dựng bộ giải mã hành động có cấu trúc vật lý. Đầu ra
chuẩn hóa của mạng nơ-ron không được ánh xạ trực tiếp lên toàn bộ miền vật lý. Thay vào
đó, tỷ lệ công suất luồng chung được đưa vào miền từ $0,70$ đến $0,99$ của tổng công suất;
phần công suất riêng còn lại được chia bằng hàm softmax có hệ số độ sắc bằng $10$; tỷ lệ
phân bổ tốc độ chung được chiếu lên đơn hình; hệ số chia năng lượng được ánh xạ tuyến tính
để tự thỏa ràng buộc bảo toàn; và pha của hai phía được học như phần hiệu chỉnh trong biên
$\pm 0,25\pi$ quanh một pha tham chiếu phụ thuộc kênh. Nhờ cách tổ chức này, mọi hành động
do mạng sinh ra đều khả thi về mặt vật lý mà không cần bước chiếu hay hạng phạt bổ sung.

Ba thuật toán DDPG, TD3 và PPO được huấn luyện trên cùng bộ giải mã, cùng ngân hàng kịch
bản kênh và cùng giao thức lựa chọn mô hình. Kết quả thực nghiệm trên $1.000$ kịch bản
kiểm thử cho thấy ba điểm chính.

Thứ nhất, không có bằng chứng rằng một thuật toán DRL luôn vượt trội. Tại $N = 16$, PPO
kém rõ so với TD3 và DDPG, nhưng từ $N = 64$ trở lên ba thuật toán đạt chất lượng rất gần
nhau, cùng nằm trong vùng $19$ đến $21$ bit/s/Hz. Sự hội tụ này ủng hộ nhận định rằng khi
biểu diễn hành động đã mã hóa phần lớn cấu trúc khó của bài toán, khác biệt giữa các cơ
chế cập nhật chính sách giảm đáng kể.

Thứ hai, nhóm tối ưu lặp đạt tổng tốc độ cao hơn. AO--Grid cao hơn toàn bộ nhóm DRL ở mọi
kích thước, với khoảng cách tại $N = 128$ vào khoảng $1,37$ đến $1,47$ bit/s/Hz. Điều này
phù hợp với việc AO tìm kiếm trên miền vật lý rộng hơn và dành nhiều phép đánh giá hàm mục
tiêu hơn cho từng CSI.

Thứ ba, chênh lệch về độ trễ lớn hơn nhiều so với chênh lệch về chất lượng. Ba thuật toán
DRL giữ độ trễ dưới một mili giây trên toàn bộ miền khảo sát, trong khi AO--Grid cần từ
$189$ ms đến hơn $1,2$ giây và AO--SCA cần từ $0,83$ đến $3,61$ giây. Khoảng cách ba đến
bốn bậc độ lớn về thời gian đổi lấy khoảng bảy phần trăm về tổng tốc độ chính là đóng góp
thực nghiệm trung tâm của luận văn.

Về QoS, nhóm DRL đạt xác suất tất cả bốn người dùng cùng đạt ngưỡng trên $98,9\%$ từ $N =
32$ trở lên. Kết quả này cho thấy việc ưu tiên tính khả thi QoS trong bước chọn điểm lưu
mô hình đã phát huy tác dụng, và tỷ lệ công suất luồng chung cao thực sự đóng vai trò mức
nền dịch vụ như phân tích đã dự đoán.

Tổng hợp lại, luận văn không đề xuất một thuật toán duy nhất cho mọi tình huống mà cung
cấp một bản đồ lựa chọn theo yêu cầu hệ thống. Khi CSI thay đổi chậm và có đủ ngân sách
tính toán, các bộ giải lặp cho chất lượng cao hơn. Khi kênh biến đổi nhanh và bộ điều
khiển phải quyết định dưới một mili giây, một chính sách DRL có biểu diễn hành động dựa
trên cấu trúc vật lý là lựa chọn thực tế hơn.

\section*{2. Hướng phát triển}
\addcontentsline{toc}{section}{2. Hướng phát triển}

Một hướng phát triển trực tiếp là xây dựng pha tham chiếu riêng cho phía truyền và phía
phản xạ. Bộ giải mã hiện tại dùng cùng một pha tham chiếu có trọng số rồi cho hai phía học
hai phần hiệu chỉnh độc lập. Trong tương lai, có thể tính hai pha tham chiếu từ riêng các
người dùng nằm ở mỗi phía để phản ánh tốt hơn cấu trúc STAR--RIS.

Hướng thứ hai là thay các hằng số cấu trúc cố định bằng tham số có thể thích nghi. Chẳng
hạn, cận dưới và cận trên của tỷ lệ công suất luồng chung có thể được điều chỉnh theo tỷ
số tín hiệu trên tạp âm, yêu cầu QoS hoặc đặc điểm kênh; hệ số độ sắc của softmax cũng có
thể được học hoặc thay đổi theo quá trình huấn luyện. Cách này có thể giữ lợi ích của biểu
diễn có cấu trúc nhưng giảm nguy cơ thu hẹp miền hành động quá mức.

Hướng thứ ba là nghiên cứu cơ chế tạo phân bổ thưa chính xác hơn. Hàm softmax hiện tại chỉ
tạo các trọng số rất gần không chứ không bằng không tuyệt đối. Các phép biến đổi có khả
năng tạo nghiệm thưa hoặc cơ chế chọn người dùng rời rạc kết hợp tối ưu liên tục có thể
phù hợp với trực giác một luồng riêng chiếm ưu thế của hệ SISO.

Hướng thứ tư là mở rộng mô hình sang hệ nhiều anten. Khi trạm gốc có nhiều anten, khả năng
tạo búp sóng theo không gian làm thay đổi hoàn toàn quan hệ nhiễu giữa các luồng riêng.
Khi đó, kết luận về xu hướng tập trung công suất riêng của hệ SISO không thể chuyển nguyên
vẹn sang hệ nhiều đầu vào, một đầu ra (Multiple-Input Single-Output, MISO) hoặc nhiều đầu
vào, nhiều đầu ra (Multiple-Input Multiple-Output, MIMO); bộ giải mã cần được thiết kế lại
để kết hợp tiền mã hóa không gian, RSMA và STAR--RIS.

Hướng thứ năm là xét CSI không hoàn hảo, lượng tử hóa pha, tổn hao phần tử, quan hệ ghép
pha truyền và phản xạ cùng chuyển động người dùng. Đây là bước cần thiết để đánh giá độ
bền của phương pháp khi chuyển từ mô hình lý tưởng sang điều kiện gần thực tế hơn.

Hướng thứ sáu là kết hợp DRL và AO. Một chính sách DRL có thể tạo nghiệm ban đầu rất
nhanh, sau đó AO thực hiện một số ít bước tinh chỉnh khi thời gian cho phép. Cách kết hợp
này có tiềm năng nằm giữa hai cực: nhanh như chính sách học khi cần quyết định tức thời,
nhưng có khả năng tiến gần chất lượng AO khi hệ thống có thêm ngân sách tính toán.

Cuối cùng, cần mở rộng đánh giá sang số lượng hạt giống lớn hơn, nhiều mô hình kênh hơn và
nền tảng phần cứng thực hoặc mô phỏng phần cứng trong vòng lặp. Khi đó, có thể kiểm tra
không chỉ tổng tốc độ và QoS mà còn mức tiêu thụ năng lượng tính toán, độ trễ đầu cuối, độ
bền trước sai số CSI và khả năng tổng quát hóa khi cấu hình hệ thống thay đổi.
